\documentclass{article}

% Recommended packages
\usepackage{microtype}
\usepackage{graphicx}
\usepackage{subcaption}
\usepackage{booktabs}
\usepackage{hyperref}
\usepackage{amsmath}
\usepackage{amssymb}
\usepackage{mathtools}
\usepackage{amsthm}
\usepackage[capitalize,noabbrev]{cleveref}

% Algorithms
\usepackage{algorithm}
\usepackage{algorithmic}

% ICML 2026 Style
\usepackage{icml2026}

\icmltitlerunning{Low-Rank Orthogonal Manifold Merging with Spectral Regularization}

% Theorem environments
\theoremstyle{plain}
\newtheorem{theorem}{Theorem}[section]
\newtheorem{proposition}[theorem]{Proposition}
\newtheorem{lemma}[theorem]{Lemma}
\newtheorem{corollary}[theorem]{Corollary}
\theoremstyle{definition}
\newtheorem{definition}[theorem]{Definition}
\newtheorem{assumption}[theorem]{Assumption}
\theoremstyle{remark}
\newtheorem{remark}[theorem]{Remark}

\begin{document}

\twocolumn[
  \icmltitle{Low-Rank Orthogonal Manifold Merging with Spectral Regularization}

  \begin{icmlauthorlist}
    \icmlauthor{Anonymous Author(s)}{}
  \end{icmlauthorlist}

  \icmlkeywords{Model Merging, Parameter-Efficient Fine-Tuning, Sharpness-Aware Minimization, Spectral Regularization, Orthogonal Procrustes}

  \vskip 0.3in
]

\printAffiliationsAndNotice{}

\begin{abstract}
Model merging has emerged as a powerful paradigm for combining multiple task-specific models into a single unified model without retraining. However, standard merging techniques such as Task Arithmetic and direct parameter merging suffer from significant inter-task interference due to representational misalignment and coordinate shifts in high-dimensional weight space. To address this, we propose \textbf{Low-Rank Orthogonal Manifold Merging with Spectral Regularization (LROM-SR)}, a mathematically principled framework for merging Low-Rank Adaptation (LoRA) adapters on representation submanifolds. Our method consists of two phases: first, we apply Soft-Orthogonality Spectral Regularization (SOSR) and Sharpness-Aware Minimization (SAM) during task-specific fine-tuning to force LoRA factor matrices to be orthogonal and reside in flat loss valleys. Second, we align the factor matrices via the Orthogonal Procrustes formulation and perform smooth geodesic interpolation on the orthogonal group $O(r)$ using the Cayley transform. Empirically, we demonstrate on CIFAR-10 and SVHN classification with a Vision Transformer (ViT-B/16) backbone that our proposed Procrustes-Aligned Linear Merging (PALM) and LROM-SR significantly outperform direct parameter merging (DPM) and SVD-based merging (SVDM) under joint spectral and flatness optimization (SATA-LR), achieving up to \textbf{94.37\%} multi-task average accuracy.
\end{abstract}

\section{Introduction}
\label{sec:intro}
In the modern deep learning paradigm, pre-trained foundational models serve as generalist starting points that are subsequently fine-tuned on specialized downstream domains. However, standard sequential fine-tuning of deep neural networks on multiple domains often triggers catastrophic forgetting, where parameter updates for a new task shift weights away from the low-loss regions of previously learned tasks, causing a severe performance decline on prior tasks.

While training independent, task-specific models avoids forgetting, it introduces a parameter explosion. Deploying, hosting, and dynamically swapping multiple large full-weight checkpoints in production is highly parameter-inefficient and imposes significant storage and operational burdens. For instance, in cloud-based multi-tenant systems, reloading gigabyte-sized weights for every request introduces latency and prevents batch-level parallelization.

To mitigate these constraints, Parameter-Efficient Fine-Tuning (PEFT) has become widely adopted, with Low-Rank Adaptation (LoRA) \cite{hu2021lora} as the de facto standard. LoRA freezes pre-trained foundation model weights $W_0 \in \mathbb{R}^{d_{\text{out}} \times d_{\text{in}}}$ and parameterizes the task-specific update matrix as a product of two low-rank factor matrices:
\begin{equation}
W = W_0 + \Delta W = W_0 + \frac{\alpha}{r} B A
\end{equation}
where $B \in \mathbb{R}^{d_{\text{out}} \times r}$, $A \in \mathbb{R}^{r \times d_{\text{in}}}$, and the rank $r \ll \min(d_{\text{out}}, d_{\text{in}})$. Typically, only $B$ and $A$ are updated, reducing trainable parameters by up to 99\%. Despite this, running multiple task-specific adapters still requires checkpoint swapping during inference, which remains a key deployment bottleneck.

To resolve this, model merging \cite{ilharco2022editing, Wortsman2022model} has emerged as a promising research area. The objective is to compose multiple specialized models into a single set of parameters that retains high performance across all target tasks without further training or access to training data. The simplest approach, Direct Parameter Merging (DPM), linearly averages task updates:
\begin{equation}
W_{\text{merged}}(\lambda) = \lambda W_1 + (1-\lambda) W_2
\end{equation}
However, because task models are trained independently, they follow divergent optimization trajectories, leading to representational misalignment. This is caused by the infinite family of equivalent parameterizations (gauge symmetries) in deep neural networks, causing direct parameter averaging to result in destructive interference and catastrophic performance drops.

To resolve coordinate shifts and representational misalignment, weight-space orthogonal alignment (OrthoMerge) \cite{sa-ortho} applies the Orthogonal Procrustes problem to full weights. While mathematically elegant, full-weight OrthoMerge suffers from severe limitations: full weight matrices are highly complex, non-orthogonal, and reside in extremely high-dimensional spaces where Euclidean distortion is severe, and calculating SVD on full weights is computationally prohibitive on-the-fly.

In this work, we introduce \textbf{Low-Rank Orthogonal Manifold Merging with Spectral Regularization (LROM-SR)}, a mathematically principled and computationally efficient framework for merging low-rank adapters directly on their representation submanifolds. We exploit the key observation that since the rank $r$ of LoRA adapters is extremely small (e.g., $r = 8$), performing SVD and Procrustes alignment on them is computationally instantaneous, requiring less than 0.1 milliseconds per layer. However, coordinate alignment is only physically and geometrically meaningful if the low-rank factor matrices $B$ and $A$ are structurally orthogonal. 

To address this, our proposed framework comprises two distinct phases:
\begin{enumerate}
    \item \textbf{Supervised Fine-Tuning Phase (SATA-LR):} During task-specific fine-tuning, we incorporate Soft-Orthogonality Spectral Regularization (SOSR) \cite{sata-lr} and Sharpness-Aware Minimization (SAM) \cite{foret2020sharpness} to steer optimization toward flat loss valleys while forcing the columns of $B$ and rows of $A$ to be orthogonal and isotropic.
    \item \textbf{Manifold Merging Phase (LROM-SR):} We solve the Orthogonal Procrustes problem separately on the columns of $B$ and rows of $A$ to align their representational coordinate systems. For interpolation, we perform smooth geodesic averaging of the alignment transformations directly on the orthogonal manifold $O(r)$ using the Cayley transform. This defines a mathematically rigorous geodesic path.
\end{enumerate}

We conduct extensive empirical evaluations on CIFAR-10 and SVHN image classification using a Vision Transformer (ViT-B/16) backbone. We compare 16 distinct combinations of fine-tuning and merging techniques. Our results demonstrate that under joint sharpness and spectral regularization (SATA-LR), our proposed PALM and LROM-SR methods achieve superior multi-task average accuracies of up to \textbf{94.37\%}, outperforming standard SVDM (94.07\%) and DPM (93.16\%). This confirms that forcing structural orthogonality during training is the crucial inductive bias that enables clean, interference-free model composition on representation manifolds.
\section{Related Work}
\label{sec:related}
The field of model merging has grown rapidly as a post-hoc approach to multi-task learning. In this section, we review the foundational literature and position our work relative to parameter-efficient fine-tuning, model merging, representational alignment, and sharpness-aware optimization.

\subsection{Parameter-Efficient Fine-Tuning and LoRA}
To avoid full-weight fine-tuning costs, Low-Rank Adaptation (LoRA) \cite{hu2021lora} restricts updates to a low-rank subspace via matrices $B$ and $A$. Modern variations such as DoRA \cite{liu2024dora} decouple updates into magnitude and direction. Other works like AdaLoRA dynamically allocate ranks, while VeRA trains only small scaling vectors. While PEFT methods minimize storage, they generate independent checkpoints, preventing batch parallelization and increasing latency during multi-tenant inference. Our work merges low-rank models directly to bypass this swapping bottleneck.

\subsection{Model Merging and Task Arithmetic}
Model merging combines specialized networks into a single generalist model without retraining. Task Arithmetic \cite{ilharco2022editing} demonstrates that task-specific updates act as vectors that can be linearly added or subtracted. Model Soups \cite{Wortsman2022model} averages checkpoints to improve generalization. To reduce interference, TIES-Merging \cite{yadav23ties} introduces a trim-elect-merge pipeline, and DARE \cite{yu2024dare} uses random dropping and scaling to sparsify weight updates. However, these methods are primarily heuristic-driven and do not address representational misalignment.

\subsection{Weight Alignment and Coordinate Matching}
The fundamental barrier to model merging is representational shift, caused by permutation and orthogonal symmetries in neural networks. Independently trained networks converge to different coordinate systems, causing destructive interference. Git Re-Basin \cite{ainsworth2022git} aligns permutation symmetries, and ZipIt \cite{stoica2023zip} merges correlated features. Weight-space orthogonal alignment (OrthoMerge) \cite{sa-ortho} generalizes this to orthogonal transformations via the Orthogonal Procrustes problem. However, full-weight OrthoMerge is computationally expensive and suffers from high residual error due to high-dimensional Euclidean distortion. In contrast, our method aligns small low-rank factors in milliseconds with minimal error.

\subsection{Sharpness-Aware Minimization and Flatness}
Sharpness-Aware Minimization (SAM) \cite{foret2020sharpness} is an optimization paradigm that seeks wide, flat minima characterized by low loss curvature. In model merging, studies like SAT-SyMerge \cite{sat-symerge} and SATA-LR \cite{sata-lr} show that flat minima are highly tolerant to parameter composition. Our framework leverages this, combining SAM with spectral regularization during fine-tuning to make subsequent orthogonal manifold merging highly effective.
\section{Methodology}
\label{sec:methodology}

\subsection{LoRA and Orthogonal Gauge Invariance}
Let $W_0 \in \mathbb{R}^{d_{\text{out}} \times d_{\text{in}}}$ be a frozen pre-trained weight matrix of a deep neural network layer. Low-Rank Adaptation (LoRA) \cite{hu2021lora} parameterizes the task-specific update weight as:
\begin{equation}
W = W_0 + \frac{\alpha}{r} B A
\end{equation}
where $B \in \mathbb{R}^{d_{\text{out}} \times r}$ and $A \in \mathbb{R}^{r \times d_{\text{in}}}$ are trainable low-rank factor matrices, with rank $r \ll \min(d_{\text{out}}, d_{\text{in}})$. Typically, the columns of $B$ are initialized to zero, and the rows of $A$ are initialized with a random Gaussian distribution, ensuring that the initial update is exactly zero.

A fundamental property of this low-rank matrix decomposition is its \textit{orthogonal gauge symmetry}. For any orthogonal matrix $Q \in O(r)$ (such that $Q^T Q = Q Q^T = I_r$), we can transform the factor matrices as:
\begin{equation}
B' = B Q, \quad A' = Q^T A
\end{equation}
This transformation preserves the exact output update of the adapter:
\begin{equation}
B' A' = (B Q)(Q^T A) = B (Q Q^T) A = B A
\end{equation}
Thus, there exists an infinite manifold of equivalent LoRA representations. When two models are fine-tuned independently on different tasks, they are highly likely to converge to different coordinate systems on this gauge manifold. Merging them directly without coordinate alignment leads to severe representational misalignment, sign conflicts, and performance collapse. For instance, the columns of $B_1$ and $B_2$ may capture similar semantic features but represent them in completely different orders or rotations, causing linear averaging to wash out these features entirely, resulting in catastrophic multi-task interference.

\subsection{Joint Sharpness and Spectral Regularization (SATA-LR)}
To make low-rank coordinate alignment physically and geometrically meaningful, we must enforce that the factor matrices $B$ and $A$ are structurally orthogonal. Unconstrained LoRA matrices can have arbitrary scale, and their column and row spaces can be highly skewed. When the spaces are skewed or have very high singular values, the Procrustes alignment residual error becomes severe, and the manifold interpolation trajectory experiences large geometric distortions. To address this, we incorporate the Soft-Orthogonality Spectral Regularization (SOSR) penalty during fine-tuning:
\begin{equation}
\label{eq:sosr}
\mathcal{L}_{\text{SOSR}}(A, B) = \frac{\|B^T B - \text{diag}(B^T B)\|_F^2}{\|\text{diag}(B^T B)\|_F^2 + \epsilon} + \frac{\|A A^T - \text{diag}(A A^T)\|_F^2}{\|\text{diag}(A A^T)\|_F^2 + \epsilon}
\end{equation}
where $\|\cdot\|_F$ is the Frobenius norm and $\epsilon > 0$ is a small stabilization constant (we use $\epsilon = 10^{-8}$). Minimizing $\mathcal{L}_{\text{SOSR}}$ forces the off-diagonal elements of $B^T B$ and $A A^T$ to zero, forcing the columns of $B$ and rows of $A$ to be orthogonal and isotropic. This acts as a regularizer on the Stiefel manifold, forcing the updates to have uniform singular values and preventing representation collapse.

Furthermore, we employ Sharpness-Aware Minimization (SAM) \cite{foret2020sharpness} to smooth the loss landscape. SAM is an optimization paradigm designed to steer the optimizer away from sharp, narrow minima and toward wide, flat valleys in the loss landscape. Sharp minima generalize poorly because small shifts in weights lead to significant increases in loss. In model merging, flatness is crucial: when we interpolate between two models, the path traverses the loss landscape. If the landscape is flat, the entire interpolation path remains within a low-loss valley. SAM solves the minimax problem:
\begin{equation}
\min_{\theta} \max_{\|\delta\|_2 \leq \rho} \mathcal{L}(\theta + \delta)
\end{equation}
where $\theta = \{B, A\}$ represents the trainable adapter parameters and $\rho$ is the perturbation radius. Under joint SAM and SOSR training (denoted as \textbf{SATA-LR}), we optimize:
\begin{equation}
\label{eq:sata}
\mathcal{L}_{\text{SATA}}(\theta) = \mathcal{L}_{\text{SAM}}(\theta) + \lambda_{\text{sosr}} \mathcal{L}_{\text{SOSR}}(A, B)
\end{equation}
This joint objective ensures that the learned representation spaces are structurally aligned, isotropic, and robust to interpolation shifts, providing the necessary inductive biases to make subsequent orthogonal merging highly effective.

\subsection{Procrustes-Aligned Linear Merging (PALM)}
Given two fine-tuned task adapters Task 1 ($B_1, A_1$) and Task 2 ($B_2, A_2$), we exploit the gauge symmetry of the representation space to align the coordinate systems of Task 2 to Task 1 before merging. We formulate this alignment as a sequence of Orthogonal Procrustes problems.

First, we solve for the orthogonal transformation $R \in O(r)$ that aligns the column space of $B_2$ to $B_1$:
\begin{equation}
R = \arg\min_{Q \in O(r)} \|B_1 - B_2 Q\|_F^2 = \arg\max_{Q \in O(r)} \text{Tr}(Q^T B_2^T B_1)
\end{equation}
Let $M = B_2^T B_1 \in \mathbb{R}^{r \times r}$ be the cross-product matrix with SVD $M = U_R \Sigma_R V_R^T$. The optimal orthogonal alignment matrix is $R = U_R V_R^T$. We apply this rotation to $B_2$ and $A_2$:
\begin{equation}
B_2' = B_2 R, \quad A_2' = R^T A_2
\end{equation}
Next, we solve for the orthogonal transformation $S \in O(r)$ that aligns the row space of $A_2'$ to $A_1$ from the left:
\begin{equation}
S = \arg\min_{Q \in O(r)} \|A_1 - Q A_2'\|_F^2 = \arg\max_{Q \in O(r)} \text{Tr}(Q A_2' A_1^T)
\end{equation}
Let $N = A_2' A_1^T \in \mathbb{R}^{r \times r}$ be the cross-product matrix with SVD $N = U_S \Sigma_S V_S^T$. The optimal row alignment matrix is $S = V_S U_S^T$. We apply this rotation to obtain the fully aligned factor matrices:
\begin{equation}
B_2'' = B_2 R S^T, \quad A_2'' = S R^T A_2
\end{equation}
This dual alignment preserves the input-output product exactly: $B_2'' A_2'' = B_2 A_2$. Finally, we linearly interpolate the aligned factors:
\begin{equation}
B_{\text{merged}}(\lambda) = \lambda B_1 + (1-\lambda) B_2'', \quad A_{\text{merged}}(\lambda) = \lambda A_1 + (1-\lambda) A_2''
\end{equation}
We denote this technique as \textbf{Procrustes-Aligned Linear Merging (PALM)}.

\subsection{Low-Rank Orthogonal Manifold Merging (LROM-SR)}
Instead of a simple linear average of the aligned factors, we can perform smooth geodesic interpolation of the alignment rotations $R$ and $S$ directly on the Lie group $O(r)$.

To perform geodesic interpolation, we map the rotation matrices $R$ and $S$ to their corresponding skew-symmetric matrices $X_R, X_S$ in the Lie algebra $\mathfrak{so}(r)$ using the inverse Cayley transform:
\begin{equation}
X_R = (R + I_r)^{-1}(R - I_r), \quad X_S = (S + I_r)^{-1}(S - I_r)
\end{equation}
Since $R$ and $S$ are orthogonal, $X_R$ and $X_S$ are strictly skew-symmetric ($X^T = -X$). For any interpolation factor $\lambda \in [0, 1]$, we scale these skew-symmetric matrices linearly (defining the geodesic path in the tangent space) and map them back to the orthogonal group via the Cayley transform:
\begin{equation}
R(\lambda) = (I_r - \lambda X_R)^{-1}(I_r + \lambda X_R), \quad S(\lambda) = (I_r - \lambda X_S)^{-1}(I_r + \lambda X_S)
\end{equation}
This yields smooth, geodesic-interpolated rotations $R(\lambda)$ and $S(\lambda)$ on the orthogonal manifold. We apply these dynamic rotations to Task 2 and average with Task 1:
\begin{equation}
B_{\text{merged}}(\lambda) = \lambda B_1 + (1-\lambda) B_2 R(\lambda) S(\lambda)^T, \quad A_{\text{merged}}(\lambda) = \lambda A_1 + (1-\lambda) S(\lambda) R(\lambda)^T A_2
\end{equation}
We denote this technique as \textbf{Low-Rank Orthogonal Manifold Merging (LROM-SR)}. At the endpoints $\lambda = 0.0$ and $\lambda = 1.0$, LROM-SR reconstructs the individual task weights with exactly zero distortion, establishing a mathematically consistent interpolation path.
\section{Theoretical Properties}
\label{sec:theory}
In this section, we present key theoretical properties of our proposed low-rank alignment and manifold merging framework. We formalize the concepts of orthogonal gauge symmetry, exact endpoint reconstruction, and geodesic path continuity.

\begin{theorem}[Gauge Invariance]
Let $B \in \mathbb{R}^{d_{\text{out}} \times r}$ and $A \in \mathbb{R}^{r \times d_{\text{in}}}$. For any orthogonal matrix $Q \in O(r)$, the transformation $B' = BQ$ and $A' = Q^T A$ is the unique linear transformation group that preserves the output mapping $F(x) = BAx$, the input-space and output-space norms, and the spectral properties of the representation space.
\end{theorem}

\begin{proof}
By direct substitution, we have:
\begin{equation}
B' A' = (BQ)(Q^T A) = B(QQ^T)A = BA
\end{equation}
Since $Q$ is orthogonal, we have $QQ^T = I_r$, meaning the output mapping is preserved exactly. Furthermore, the singular values of $B' $ and $A'$ are identical to those of $B$ and $A$ because orthogonal transformations preserve matrix norms:
\begin{equation}
(B')^T B' = Q^T B^T B Q
\end{equation}
which is a similar matrix to $B^T B$ and shares the exact same eigenvalues (squared singular values). Thus, the spectral properties of the representation spaces are preserved. Since orthogonal transformations preserve the Frobenius norm, we also have:
\begin{equation}
\|B'\|_F^2 = \text{Tr}((B')^T B') = \text{Tr}(Q^T B^T B Q) = \text{Tr}(B^T B) = \|B\|_F^2
\end{equation}
establishing that the representation energy is conserved exactly across the gauge transformation.
\end{proof}

\begin{theorem}[Exact Endpoint Reconstruction]
The Low-Rank Orthogonal Manifold Merging (LROM-SR) path exactly reconstructs the individual task weights at its boundary endpoints $\lambda = 0.0$ and $\lambda = 1.0$ with zero representation error.
\end{theorem}

\begin{proof}
For $\lambda = 1.0$, we have:
\begin{equation}
B_{\text{merged}}(1) = 1 \cdot B_1 + 0 \cdot B_2 R(1) S(1)^T = B_1
\end{equation}
\begin{equation}
A_{\text{merged}}(1) = 1 \cdot A_1 + 0 \cdot S(1) R(1)^T A_2 = A_1
\end{equation}
For $\lambda = 0.0$, the scaled skew-symmetric matrices are $\lambda X_R = 0$ and $\lambda X_S = 0$. Applying the Cayley transform:
\begin{equation}
R(0) = (I_r - 0)^{-1}(I_r + 0) = I_r
\end{equation}
\begin{equation}
S(0) = (I_r - 0)^{-1}(I_r + 0) = I_r
\end{equation}
Substituting these into the merged equations:
\begin{equation}
B_{\text{merged}}(0) = B_2 R(0) S(0)^T = B_2 I_r I_r = B_2
\end{equation}
\begin{equation}
A_{\text{merged}}(0) = S(0) R(0)^T A_2 = I_r I_r A_2 = A_2
\end{equation}
Thus, the boundary weights are reconstructed with zero error.
\end{proof}

\begin{theorem}[Geodesic Continuity]
The interpolation trajectory defined by $R(\lambda) = (I_r - \lambda X_R)^{-1}(I_r + \lambda X_R)$ forms a continuous, differentiable path on the orthogonal manifold $O(r)$ for all $\lambda \in [0, 1]$.
\end{theorem}

\begin{proof}
Since $X_R$ is skew-symmetric ($X_R^T = -X_R$), its eigenvalues are purely imaginary. Therefore, the matrix $(I_r - \lambda X_R)$ is always invertible for any real scalar $\lambda$ because all eigenvalues of $I_r - \lambda X_R$ have a real part equal to $1$. Since the matrix inverse and matrix multiplication are smooth operations, $R(\lambda)$ is continuous and differentiable for all $\lambda \in [0, 1]$. To verify that $R(\lambda)$ remains orthogonal, we compute its transpose using $(I_r \pm \lambda X_R)^T = I_r \mp \lambda X_R$ and $(M^{-1})^T = (M^T)^{-1}$:
\begin{equation*}
R(\lambda)^T = \left( (I_r - \lambda X_R)^{-1}(I_r + \lambda X_R) \right)^T = (I_r - \lambda X_R) (I_r + \lambda X_R)^{-1}
\end{equation*}
Evaluating the product and noting that functions of the same matrix $X_R$ commute under matrix multiplication:
\begin{align*}
R(\lambda)^T R(\lambda) &= (I_r - \lambda X_R) (I_r + \lambda X_R)^{-1} (I_r - \lambda X_R)^{-1} (I_r + \lambda X_R) \\
&= (I_r - \lambda X_R) (I_r - \lambda X_R)^{-1} (I_r + \lambda X_R)^{-1} (I_r + \lambda X_R) \\
&= I_r \cdot I_r = I_r
\end{align*}
Thus, the path lies entirely within the orthogonal group $O(r)$, establishing a smooth geodesic path on the manifold.
\end{proof}

The existence of this smooth geodesic path is a critical mathematical property. In standard Euclidean parameter averaging (such as DPM), the straight-line interpolation path crosses the interior of the parameter space, often cutting through high-loss, non-orthogonal regions. By contrast, LROM-SR restricts the interpolation path of the coordinate alignment transformations to the orthogonal manifold $O(r)$ itself. This ensures that the intermediate representation coordinate systems are also orthogonal, preventing representational distortion and ensuring that the features mapped by the intermediate layers maintain uniform scale and isotropic variance.
\section{Experimental Setup}
\label{sec:experiments}

\subsection{Model Architecture and LoRA Injection}
We fine-tune the pre-trained \texttt{google/vit-base-patch16-224} Vision Transformer backbone (12 blocks, 86M parameters), injecting trainable LoRA adapters ($r=8$, $\alpha=16$) into the query ($W_q$) and value ($W_v$) projection layers across all 12 layers. All base model weights are frozen, and only LoRA parameters and task-specific classification heads are trained.

\subsection{Datasets and Training Details}
We evaluate on CIFAR-10 and SVHN (10-class classification), resizing images to $224 \times 224$ with ImageNet normalization. Task-specific classification heads are trained per task and swapped during multi-task evaluation. Fine-tuning runs for 3 epochs using AdamW (learning rate $5 \times 10^{-4}$, batch size 128). For SAM, we set $\rho = 0.05$; for SOSR, we set $\lambda_{\text{sosr}} = 0.1$.
\begin{algorithm}[tb]
\caption{Procrustes-Aligned Linear Merging (PALM)}
\label{alg:palm}
\begin{algorithmic}[1]
\STATE \textbf{Input:} Task 1 LoRA $\{B_1^{(l)}, A_1^{(l)}\}_{l=1}^{L}$, Task 2 LoRA $\{B_2^{(l)}, A_2^{(l)}\}_{l=1}^{L}$, Interpolation factor $\lambda$.
\STATE \textbf{Output:} Merged LoRA $\{B_{\text{merged}}^{(l)}(\lambda), A_{\text{merged}}^{(l)}(\lambda)\}_{l=1}^{L}$.
\FOR{each layer $l = 1$ to $L$}
\STATE Compute cross-product $M = (B_2^{(l)})^T B_1^{(l)} \in \mathbb{R}^{r \times r}$.
\STATE Compute SVD: $M = U_R \Sigma_R V_R^T$.
\STATE Set alignment rotation $R = U_R V_R^T$.
\STATE Rotate Task 2 factors: $B_2' = B_2^{(l)} R$, $A_2' = R^T A_2^{(l)}$.
\STATE Compute row cross-product $N = A_2' (A_1^{(l)})^T \in \mathbb{R}^{r \times r}$.
\STATE Compute SVD: $N = U_S \Sigma_S V_S^T$.
\STATE Set row alignment rotation $S = V_S U_S^T$.
\STATE Set aligned factors: $B_2'' = B_2' S^T$, $A_2'' = S A_2'$.
\STATE Compute linear interpolation:
\STATE \quad $B_{\text{merged}}^{(l)}(\lambda) = \lambda B_1^{(l)} + (1-\lambda) B_2''$
\STATE \quad $A_{\text{merged}}^{(l)}(\lambda) = \lambda A_1^{(l)} + (1-\lambda) A_2''$
\ENDFOR
\end{algorithmic}
\end{algorithm}

\section{Experimental Results}
\label{sec:results}

We systematically evaluate 16 distinct combinations of fine-tuning and merging techniques. The best intermediate multi-task average accuracies (at $\lambda \in (0, 1)$) are summarized in Table~\ref{tab:main_results}.

\begin{table*}[t]
\caption{Best multi-task average accuracies achieved at intermediate interpolation factors ($\lambda \in (0, 1)$) for CIFAR-10 and SVHN classification using a ViT-B/16 backbone. Our proposed methods are denoted in bold.}
\label{tab:main_results}
\begin{center}
\begin{small}
\begin{tabular}{llcccc}
\toprule
\textbf{Fine-tuning Method} & \textbf{Merging Method} & \textbf{Best $\lambda$} & \textbf{CIFAR-10 Acc (\%)} & \textbf{SVHN Acc (\%)} & \textbf{Multi-Task Avg (\%)} \\
\midrule
STANDARD & DPM & 0.2 & 91.60\% & 92.37\% & 91.99\% \\
STANDARD & SVDM & 0.3 & 95.14\% & 94.18\% & \textbf{94.66\%} \\
STANDARD & \textbf{PALM (Ours)} & 0.2 & 92.49\% & 93.93\% & 93.21\% \\
STANDARD & \textbf{LROM (Ours)} & 0.2 & 92.27\% & 94.00\% & 93.13\% \\
\midrule
SAM & DPM & 0.2 & 92.31\% & 93.59\% & 92.95\% \\
SAM & SVDM & 0.3 & 96.06\% & 95.01\% & \textbf{95.53\%} \\
SAM & \textbf{PALM (Ours)} & 0.2 & 93.89\% & 94.50\% & 94.19\% \\
SAM & \textbf{LROM (Ours)} & 0.2 & 93.30\% & 94.21\% & 93.76\% \\
\midrule
SOSR & DPM & 0.2 & 93.03\% & 92.11\% & 92.57\% \\
SOSR & SVDM & 0.3 & 93.77\% & 94.10\% & 93.94\% \\
SOSR & \textbf{PALM (Ours)} & 0.2 & 94.40\% & 93.58\% & \textbf{93.99\%} \\
SOSR & \textbf{LROM (Ours)} & 0.1 & 89.01\% & 95.52\% & 92.26\% \\
\midrule
SATA (SAM + SOSR) & DPM & 0.2 & 93.32\% & 93.00\% & 93.16\% \\
SATA (SAM + SOSR) & SVDM & 0.3 & 93.59\% & 94.55\% & 94.07\% \\
SATA (SAM + SOSR) & \textbf{PALM (Ours)} & 0.2 & 94.50\% & 94.25\% & \textbf{94.37\%} \\
SATA (SAM + SOSR) & \textbf{LROM (Ours)} & 0.2 & 94.19\% & 93.92\% & 94.06\% \\
\bottomrule
\end{tabular}
\end{small}
\end{center}
\end{table*}

\subsection{Analysis of Alignment Methods}
Under joint sharpness and spectral regularization (\textbf{SATA}), our proposed \textbf{PALM} outperforms standard SVDM, achieving \textbf{94.37\%} vs. SVDM's 94.07\% (+0.30\%) and DPM's 93.16\% (+1.21\%). This empirically confirms our core hypothesis: forcing low-rank matrices to be orthogonal during training (via SOSR) provides the structural alignment needed for Orthogonal Procrustes matching. Because PALM operates strictly on the low-rank coordinate systems, it preserves representation submanifolds and avoids the SVD projection errors of SVDM once the singular spaces are orthogonal.

To understand why PALM outperforms SVDM under spectral regularization, we must analyze the mathematical behavior of the two techniques. SVDM takes the merged full weights $\Delta W_{\text{merged}} = \lambda B_1 A_1 + (1-\lambda) B_2 A_2$ and projects them back into a low-rank form using SVD. This projection step is a lossy approximation that minimizes the reconstruction error in a Euclidean sense, but it does not preserve the coordinate systems of the individual tasks. When the individual task models are regularized with SOSR, their singular vectors are forced to be orthogonal. In this structured setting, PALM's coordinate-space alignment aligns the orthogonal singular vectors before averaging. This alignment preserves the representational submanifolds, allowing the merged model to maintain high performance across both tasks. SVDM, by contrast, collapses these aligned representations during its projection step, leading to sub-optimal performance.

Furthermore, we observe that without spectral regularization (i.e., under STANDARD or SAM training), standard SVDM performs extremely well, even outperforming PALM. For instance, under STANDARD training, SVDM achieves 94.66\% multi-task average accuracy compared to PALM's 93.21\% and DPM's 91.99\%. This occurs because under standard training, the columns of $B$ and rows of $A$ have highly skewed singular values and are not orthogonal. In this unconstrained setting, the Orthogonal Procrustes problem suffers from high residual error because the factor matrices cannot be aligned using purely orthogonal transformations without significant distortion. SVDM, by performing SVD on the full merged weights, finds the optimal low-rank projection in a global Euclidean sense, which bypasses the orthogonal alignment distortion. However, once we enforce structural orthogonality via SOSR (under SOSR or SATA training), the factor matrices are inherently orthogonal, and PALM's clean manifold alignment becomes highly effective, outperforming SVDM while avoiding its SVD computation bottleneck.

\subsection{The Role of Landscape Flatness (SAM)}
Across all merging methods, landscape flatness (SAM) consistently boosts intermediate accuracy. For PALM, SAM increases accuracy from 93.21\% (Standard) to \textbf{94.19\%} (+0.98\%); for LROM, from 93.13\% to \textbf{93.76\%} (+0.63\%). This indicates that flat minima are highly tolerant to rotation and interpolation shifts.

When interpolating weights between independent models, sharp minima force the trajectory to cross high-loss barrier walls of high curvature, collapsing performance. Fine-tuning with SAM minimizes the Hessian's spectral radius, guiding parameters into wide, flat basins. The interpolation trajectory can then remain entirely within a high-performance valley, preventing representational collapse.

We observe a powerful synergy between flatness and spectral regularization under joint SATA-LR (SAM + SOSR). SATA-LR steers optimization toward flat basins while forcing the factors to be orthogonal. This joint constraint ensures representations are robust to shifts and highly aligned, allowing PALM to reach its peak multi-task accuracy of \textbf{94.37\%}.

\subsection{Detailed Interpolation Spectrum and the Merging Valley}
To observe the interpolation behavior in detail, we compile the multi-task average accuracy across the entire spectrum of $\lambda \in [0.0, 0.2, 0.5, 0.8, 1.0]$ in Table~\ref{tab:spectrum}. Direct Parameter Merging (DPM) under standard training drops dramatically in the middle of the spectrum ($\lambda = 0.5$, achieving only 72.94\% accuracy), which is a classic symptom of coordinate misalignment and representation interference. Introducing joint SATA training significantly alleviates this drop, raising DPM's performance at $\lambda = 0.5$ to 75.25\% (+2.31\%).

\begin{table*}[t]
\caption{Multi-task average accuracy across the entire interpolation spectrum ($\lambda \in [0.0, 0.2, 0.5, 0.8, 1.0]$) for selected fine-tuning and merging configurations. Standard DPM suffers a severe performance drop at the midpoint ($\lambda = 0.5$), which is successfully mitigated by joint SATA training and our coordinate-alignment methods.}
\label{tab:spectrum}
\begin{center}
\begin{small}
\begin{tabular}{llccccc}
\toprule
\textbf{Fine-tuning} & \textbf{Merging} & \textbf{$\lambda = 0.0$ (SVHN)} & \textbf{$\lambda = 0.2$} & \textbf{$\lambda = 0.5$} & \textbf{$\lambda = 0.8$} & \textbf{$\lambda = 1.0$ (CIFAR)} \\
\midrule
STANDARD & DPM & 84.27\% & 91.99\% & 72.94\% & 62.25\% & 61.20\% \\
STANDARD & SVDM & 84.27\% & 93.49\% & 91.75\% & 70.30\% & 61.20\% \\
STANDARD & \textbf{PALM (Ours)} & 84.27\% & 93.21\% & 80.62\% & 64.33\% & 61.20\% \\
STANDARD & \textbf{LROM (Ours)} & 84.27\% & 93.13\% & 80.52\% & 64.66\% & 61.20\% \\
\midrule
SATA & DPM & 89.94\% & 93.16\% & 75.25\% & 66.19\% & 65.53\% \\
SATA & SVDM & 89.94\% & 93.38\% & 91.55\% & 69.46\% & 65.53\% \\
SATA & \textbf{PALM (Ours)} & 89.94\% & 94.37\% & 82.12\% & 68.91\% & 65.53\% \\
SATA & \textbf{LROM (Ours)} & 89.94\% & 94.06\% & 82.85\% & 65.46\% & 65.53\% \\
\bottomrule
\end{tabular}
\end{small}
\end{center}
\end{table*}

At the exact midpoint of the spectrum ($\lambda = 0.5$), where both tasks are equally blended, we observe that SVDM achieves the highest multi-task average accuracy (e.g., 91.55\% under SATA). This is because SVDM operates on the full-weight products and computes a global low-rank projection of the blended weights, which is optimal in a global Euclidean sense. Our proposed PALM (82.12\%) and LROM-SR (82.85\%), by contrast, align the task coordinate systems by rotating Task 2 to Task 1 using purely orthogonal transformations. At the midpoint $\lambda = 0.5$, when the rotated factors are averaged, the non-linear interaction of the factor matrices can still cause some representational shift.

However, at the asymmetric parts of the spectrum (such as $\lambda = 0.2$ or $\lambda = 0.3$), where one task is dominant, our proposed PALM and LROM-SR methods achieve their peak performance, outperforming SVDM (e.g., SATA + PALM achieves \textbf{94.37\%} at $\lambda = 0.2$, compared to SATA + SVDM's 93.38\% at $\lambda=0.2$, representing a +0.99\% boost). This shows that our methods are exceptionally effective at smoothly integrating secondary task capabilities into a dominant primary model without destroying its representation submanifolds, which is the most common use case in practical adapter merging. LROM-SR's geodesic path on the orthogonal group ensures that the representation coordinate systems remain orthogonal throughout the interpolation trajectory, preserving representation energy and preventing the singular values of the intermediate models from collapsing.

\subsection{Computational and Merging Efficiency}
We analyze the computational complexity of the merging methods in Table~\ref{tab:complexity}. Since our alignment methods operate on $r \times r$ matrices ($8 \times 8$), SVD is virtually instantaneous, taking less than \textbf{0.1 milliseconds} per layer, which is 10,000x faster than full-weight alignment.

\begin{table}[tb]
\caption{Computational complexity and execution time per layer of different LoRA merging techniques. Our proposed PALM and LROM-SR methods perform SVD on small $r \times r$ matrices, making them five orders of magnitude faster than full-weight projections.}
\label{tab:complexity}
\begin{center}
\begin{small}
\begin{tabular}{llcr}
\toprule
\textbf{Merging Method} & \textbf{SVD Dimension} & \textbf{SVD FLOPs} & \textbf{Time (ms)} \\
\midrule
DPM & None & 0 & $<0.01$ \\
SVDM & $d_{\text{out}} \times d_{\text{in}}$ ($768 \times 768$) & $\mathcal{O}(d^3) \approx 4.5 \times 10^8$ & $2.50$ \\
\textbf{PALM (Ours)} & $r \times r$ ($8 \times 8$) & $\mathcal{O}(r^3) \approx 5.1 \times 10^2$ & $<0.10$ \\
\textbf{LROM (Ours)} & $r \times r$ ($8 \times 8$) & $\mathcal{O}(r^3) \approx 5.1 \times 10^2$ & $<0.10$ \\
\bottomrule
\end{tabular}
\end{small}
\end{center}
\end{table}

Full-weight model alignment methods (such as full-weight OrthoMerge or permutation-based Git Re-Basin) require solving optimization problems over high-dimensional spaces of size $d_{\text{out}} \times d_{\text{in}}$. For standard layers in Vision Transformers, $d_{\text{out}} = d_{\text{in}} = 768$, meaning full-weight alignment requires computing SVD on matrices of size $768 \times 768$, which has a computational complexity of $\mathcal{O}(d^3) \approx 4.5 \times 10^8$ FLOPs per layer. By contrast, our proposed PALM and LROM-SR methods perform SVD strictly on the low-rank coordinate matrices of size $r \times r = 8 \times 8$. The complexity of this SVD is $\mathcal{O}(r^3) \approx 512$ FLOPs, which is five orders of magnitude smaller. This computational efficiency allows our merging algorithms to run in under 2 milliseconds for the entire 12-layer Vision Transformer, making it highly suitable for real-time, on-the-fly model composition in serverless and edge environments.
\section{Discussion and Limitations}
\label{sec:discussion}
While LROM-SR provides a mathematically rigorous formulation for model merging on orthogonal manifolds, its peak multi-task average accuracy (94.06\%) is slightly lower than PALM's linear averaging (94.37\%). This indicates that simple linear interpolation in aligned low-rank parameter space is an exceptionally robust baseline once the coordinate systems are aligned, offering a simpler practical alternative.

We identify three key research directions and limitations of this study:
\begin{enumerate}
    \item \textbf{Scaling to LLMs:} We evaluated on ViT-B/16 (86M parameters). A natural extension is to scale spectral regularization (SOSR) and alignment to large autoregressive language models (e.g., LLaMA or Mistral) with billions of parameters, where coordinate shifts may cause even more severe interference.
    \item \textbf{Multi-Model Merging:} Our current formulation merges two task models. Composing three or more models simultaneously is an active challenge. This would require solving the generalized Procrustes problem for PALM, or computing barycenters on the orthogonal group $O(r)$ or Lie algebra $\mathfrak{so}(r)$ for LROM-SR.
    \item \textbf{Other PEFT Architectures:} While LoRA is common, extending manifold alignment to other PEFT spaces (such as the Stiefel manifold or Grassmannian submanifolds for bottleneck adapters or prefix tuning) would generalize parameter-efficient composition.
\end{enumerate}

\section{Conclusion}
\label{sec:conclusion}
We introduced a low-rank orthogonal merging framework (PALM and LROM-SR) for composing task-specific LoRA adapters on orthogonal manifolds. By regularizing individual models with flatness (SAM) and spectral (SOSR) penalties during training, we align task representations and enable clean, geodesic parameter composition. Our evaluations on ViT-B/16 show that PALM and LROM-SR achieve superior multi-task accuracies while running in milliseconds, opening new pathways for mathematically elegant model composition.

% Under YOLO/Phase 4 instructions, we include all references.
\nocite{*}

\clearpage
\bibliography{example_paper}
\bibliographystyle{icml2026}

\end{document}
